\section{Results}
\label{sec:results}

% Every number below is a macro: \w... from tables/v5_numbers.tex and \v... from
% tables/v4_numbers.tex (analysis/full_eval.py --campaign v5 / v4), \dg... from
% tables/diag_numbers.tex (analysis/diagnosis_macros.py), \sim... from
% tables/sim_numbers.tex (analysis/constraint_sim.py), \s... from
% tables/scale_numbers.tex (analysis/reward_scale_analysis.py). Do not type
% numbers by hand; re-run `make results`.

The main campaign (v5) trains all four arms with Reward~II, the floor on
expected content (Eq.~\ref{eq:reward-exp}). The previous campaign (v4) trained
the same four arms, with the same code, hyperparameters and budget, under
Reward~I, the floor on every sample (Eq.~\ref{eq:reward}). The two campaigns
differ in the reward and in how prompts are decoded at evaluation
(Section~\ref{sec:setup-exp}). Sections~\ref{sec:res-frontier}--\ref{sec:res-control}
report v5. Section~\ref{sec:res-diagnosis} uses v4 to show what goes wrong
under Reward~I, and why.

\subsection{The content--sentiment tradeoff curve}
\label{sec:res-frontier}

\begin{figure}[t]
\centering
\includegraphics[width=0.62\textwidth]{frontier_v5.pdf}
\caption{\textbf{Sentiment--content tradeoff curves} of the four policies
trained with Reward~II, on the $\wNsent$-sentence test split (step
$12{,}000$). Each marker is one content floor
$\tau = 100\lmb \in \{0, 5, \dots, 95\}$, scored by expected content and
expected sentiment: prompts are sampled from the policy, and each is scored by
$N=50$ task-model samples, averaged over sentences and five evaluation seeds.
Up and to the right is better on both objectives.
Figure~\ref{fig:frontier-panels} labels each curve by floor.}
\label{fig:frontier}
\end{figure}

\begin{figure}[!tp]
\centering
\includegraphics[width=0.86\textwidth]{frontier_panels_v5.pdf}
\caption{\textbf{Each policy's tradeoff curve, labelled by content floor}
(Reward~II). One panel per method, with the other three in grey. Labels give
the floor $\tau$ for every multiple of $10$. Floors whose points coincide share
one range label. Under Reward~II this happens only for the inactive floors below
the content of the most positive prompt ($\tau = 0$--$25$ for both \rrebel{}
arms). Error bars are $95\%$ bootstrap intervals over test sentences.
Compare the same panels under Reward~I in Figure~\ref{fig:frontier-v4}.}
\label{fig:frontier-panels}
\end{figure}

Figure~\ref{fig:frontier} shows the tradeoff curve each policy traces as its
conditioning input sweeps the floor from $0$ to $95$, and
Figure~\ref{fig:frontier-panels} labels every point by its floor.
Table~\ref{tab:full} summarizes each curve, and Table~\ref{tab:pairwise} gives
paired differences.

\paragraph{\rrebel{} Huber-std has the best frontier.}
Its hypervolume is $\wHuberHV$, against $\wGentHV$ for GRPO + entropy, a paired
difference of $\wGapHuberGentHV$ points
($95\%$ CI $[\wGapHuberGentHVLo, \wGapHuberGentHVHi]$). At the binding floors
($\tau = 30$--$90$) it keeps $\wGapHuberGentSbind$ more sentiment points on
average ($[\wGapHuberGentSbindLo, \wGapHuberGentSbindHi]$), and it tracks the
floors more closely too (mean deviation $\wGapHuberGentCalib$ points,
$[\wGapHuberGentCalibLo, \wGapHuberGentCalibHi]$). If a user may randomize between two of a policy's
operating points, Huber-std's curve lies above GRPO + entropy's at every
content level (smallest margin $\wDomGentMixMin$) and above the base-LM-reference
arm's by at least $\wDomGbaseMixMin$. \rrebel{} $\ell_1$-std is second
(hypervolume $\wLoneHV$, $\wGapLoneHuberHV$ against Huber-std,
$[\wGapLoneHuberHVLo, \wGapLoneHuberHVHi]$), and also clearly ahead of
GRPO + entropy ($+\wGapLoneGentHV$, $[\wGapLoneGentHVLo, \wGapLoneGentHVHi]$).
GRPO with the base-LM reference is last (hypervolume $\wGbaseHV$).

\paragraph{The GRPO policies stay diffuse.}
At $\tau = 0$, where only sentiment matters, both \rrebel{} arms reach
sentiment $\wHuberSenZero$ ($\ell_1$-std: $\wLoneSenZero$). GRPO + entropy reaches
only $\wGentSenZero$, and the base-LM-reference arm $\wGbaseSenZero$. The
reason is visible in the number of distinct prompts each policy emits across
the $\wNsent$ test sentences at a floor (Table~\ref{tab:full}): a median of
$\wHuberPrompts$ for Huber-std, against $\wGentPrompts$ for GRPO + entropy and
$\wGbasePrompts$ for the base-LM reference. Because the evaluation now samples
prompts from the policy, a policy that keeps probability on poor prompts pays
for it. Under Reward~I, which evaluated an almost-greedy prompt, GRPO + entropy
had reached $\vGentSenZero$ at $\tau = 0$.

\begin{table}[t]
\centering
\small
\setlength{\tabcolsep}{4pt}
\resizebox{\textwidth}{!}{% generated by analysis/full_eval.py -- do not edit by hand
\begin{tabular}{lccccccc}
\toprule
Arm & Hypervolume & $|\mathbb{E}[c]-\tau|$ & Shortfall & Sentiment & Distinct points & Prompts/$\tau$ & Per-sample reward \\
\midrule
R-REBEL ($\ell_1$-std) & $50.5$\,{\scriptsize$[49.7,51.3]$} & $2.8$\,{\scriptsize$[2.3,3.4]$} & $2.4$\,{\scriptsize$[1.8,2.9]$} & $44.6$\,{\scriptsize$[43.6,45.7]$} & $15/20$ & $18$ & $27.4$\,{\scriptsize$[26.7,28.1]$} \\
R-REBEL (Huber-std) & $52.7$\,{\scriptsize$[52.0,53.5]$} & $1.6$\,{\scriptsize$[1.4,1.8]$} & $1.1$\,{\scriptsize$[0.8,1.5]$} & $46.2$\,{\scriptsize$[45.1,47.3]$} & $16/20$ & $22$ & $26.2$\,{\scriptsize$[25.5,26.9]$} \\
\midrule
GRPO + entropy & $42.8$\,{\scriptsize$[42.0,43.7]$} & $2.2$\,{\scriptsize$[1.9,2.5]$} & $1.7$\,{\scriptsize$[1.3,2.2]$} & $37.6$\,{\scriptsize$[36.5,38.7]$} & $17/20$ & $88$ & $22.3$\,{\scriptsize$[21.6,23.0]$} \\
GRPO, base-LM reference & $38.9$\,{\scriptsize$[38.1,39.6]$} & $2.5$\,{\scriptsize$[2.3,2.8]$} & $2.0$\,{\scriptsize$[1.7,2.4]$} & $29.6$\,{\scriptsize$[28.7,30.6]$} & $14/20$ & $84$ & $22.9$\,{\scriptsize$[22.3,23.5]$} \\
\bottomrule
\end{tabular}
}
\caption{Summary of each arm's test-split tradeoff curve under Reward~II (step
$12{,}000$, $\wNsent$ sentences, $\wNlam$ floors). \emph{Hypervolume}: area
dominated by the curve in the $100\times100$ (content, sentiment) box, in
percent. $|\mathbb{E}[c]-\tau|$ and \emph{Shortfall}: mean distance from the
floor and mean amount by which $\mathbb{E}[c]$ falls short of it, over the
binding floors $\tau = 30$--$90$. \emph{Sentiment}: mean expected sentiment over
the same floors. \emph{Distinct points}: floors whose point does not coincide
(within $2$ points) with the previous floor's. \emph{Prompts}/$\tau$: median
number of distinct prompts across the test sentences at a floor.
\emph{Per-sample reward}: Reward~I's objective, for reference only (Reward~II
does not optimize it). Brackets are $95\%$ bootstrap intervals over sentences.}
\label{tab:full}
\end{table}

\begin{table}[t]
\centering
\small
\resizebox{\textwidth}{!}{% generated by analysis/full_eval.py -- do not edit by hand
\begin{tabular}{llccc}
\toprule
$a$ & $b$ & $\Delta$ hypervolume & $\Delta$ sentiment & $\Delta\,|\mathbb{E}[c]-\tau|$ \\
\midrule
R-REBEL (Huber-std) & GRPO + entropy & $+9.90$\,{\scriptsize$[+9.52,+10.29]$} & $+8.58$\,{\scriptsize$[+7.84,+9.28]$} & $-0.63$\,{\scriptsize$[-0.93,-0.26]$} \\
R-REBEL ($\ell_1$-std) & R-REBEL (Huber-std) & $-2.23$\,{\scriptsize$[-2.58,-1.87]$} & $-1.55$\,{\scriptsize$[-2.32,-0.75]$} & $+1.26$\,{\scriptsize$[+0.75,+1.70]$} \\
GRPO + entropy & GRPO, base-LM reference & $+3.93$\,{\scriptsize$[+3.47,+4.42]$} & $+7.98$\,{\scriptsize$[+7.41,+8.57]$} & $-0.26$\,{\scriptsize$[-0.78,+0.08]$} \\
R-REBEL ($\ell_1$-std) & GRPO + entropy & $+7.66$\,{\scriptsize$[+7.14,+8.21]$} & $+7.02$\,{\scriptsize$[+6.17,+7.85]$} & $+0.63$\,{\scriptsize$[+0.20,+1.09]$} \\
\bottomrule
\end{tabular}
}
\caption{Paired differences $a - b$ under Reward~II (step $12{,}000$), with
$95\%$ bootstrap intervals over test sentences. Each resample uses the same
sentences for both arms. $\Delta$ sentiment and
$\Delta|\mathbb{E}[c]-\tau|$ are over the binding floors $\tau=30$--$90$; a
negative $\Delta|\mathbb{E}[c]-\tau|$ means $a$ tracks the floor more closely.}
\label{tab:pairwise}
\end{table}

\subsection{Does the floor control the policy?}
\label{sec:res-control}

\begin{figure}[t]
\centering
\includegraphics[width=\textwidth]{control_v5.pdf}
\caption{\textbf{The floor controls the policy} (Reward~II, test split, step
$12{,}000$). \textbf{Left:} expected content against the floor; the dotted line
is $\mathbb{E}[c] = \tau$. \textbf{Middle:} expected sentiment.
\textbf{Right:} probability that an individual task-model sample meets the
floor. Bands are $95\%$ bootstrap intervals over sentences.}
\label{fig:control}
\end{figure}

\begin{figure}[t]
\centering
\includegraphics[width=\textwidth]{compare_v4_v5.pdf}
\caption{\textbf{Where each floor's point sits, under each reward.} Expected
content minus the floor, for every arm, at step $12{,}000$ of each campaign.
\textbf{Left:} Reward~I, the floor on every sample. \textbf{Right:} Reward~II, the
floor on expected content. Shaded: the binding floors $\tau = 30$--$90$. Under
Reward~II, $\mathbb{E}[c]$ follows $\tau$ across the whole binding range. Under
Reward~I it alternates between large surpluses and deficits.}
\label{fig:compare}
\end{figure}

A $\lmb$-conditioned policy is useful only if the floor it is given controls
what it does. Figure~\ref{fig:control} plots, against the floor, the expected
content the policy delivers, the sentiment it keeps, and the per-sample
probability that a rewrite meets the floor.

\paragraph{Expected content tracks the floor.}
For floors below the content of each arm's most positive prompt ($\tau \le 25$)
the floor is inactive and the policies correctly ignore it: they emit their
most positive prompts and over-satisfy the floor. Over the binding floors
$\tau = 30$--$90$, every arm's expected content stays close to the line
$\mathbb{E}[c] = \tau$. Huber-std's deviation ranges from $\wHuberGapMin$ to
$+\wHuberGapMax$ points (mean absolute deviation $\wHuberCalib$), $\ell_1$-std's
from $\wLoneGapMin$ to $+\wLoneGapMax$ ($\wLoneCalib$), and GRPO + entropy's from
$\wGentGapMin$ to $+\wGentGapMax$ ($\wGentCalib$). The base-LM-reference arm
tracks the floor through most of the binding range but cannot reach the highest
floors (worst deviation $\wGbaseGapMin$). Proposition~\ref{prop:penalty}
predicts a shortfall of the envelope's slope divided by $\nu$. The measured
curves' slopes over the binding floors have median $\wSlopeMed$ sentiment
points per content point, so the prediction is a fraction of a content point.
The measured deviations are a few points, and of both signs: the gap between a
converged nonparametric fixed point and a $\lmb$-conditioned network trained for
$12{,}000$ steps.

\paragraph{The comparison with Reward~I.}
Figure~\ref{fig:compare} shows the same quantity for both campaigns. Under
Reward~I, the floors' points alternate between surpluses of up to $\vGapAllMax$ content
points and deficits of up to $\vGapAllMinAbs$ over the binding floors, and the mean absolute deviation over the
binding floors is $\vHuberCalib$ for Huber-std ($\ell_1$-std: $\vLoneCalib$;
GRPO + entropy: $\vGentCalib$). Under Reward~II it is $\wHuberCalib$
($\wLoneCalib$; $\wGentCalib$). The curves also stop clustering. The number of
distinct points among the $\wNlam$ floors rises from $\vHuberDistinct$ to
$\wHuberDistinct$ for Huber-std, from $\vLoneDistinct$ to $\wLoneDistinct$ for
$\ell_1$-std, and from $\vGentDistinct$ to $\wGentDistinct$ for GRPO + entropy.
The remaining coincident points are the inactive floors. The hypervolume of the
\rrebel{} curves rises as well, from $\vHuberHV$ to $\wHuberHV$ (Huber-std) and
from $\vLoneHV$ to $\wLoneHV$ ($\ell_1$-std).

\subsection{What the policies learned}
\label{sec:res-prompts}

\begin{table}[t]
\centering
\setlength{\tabcolsep}{4pt}
% generated by analysis/full_eval.py -- do not edit by hand
\begin{tabular}{r>{\raggedright\arraybackslash\ttfamily\footnotesize}p{0.205\textwidth}>{\raggedright\arraybackslash\ttfamily\footnotesize}p{0.205\textwidth}>{\raggedright\arraybackslash\ttfamily\footnotesize}p{0.205\textwidth}>{\raggedright\arraybackslash\ttfamily\footnotesize}p{0.205\textwidth}}
\toprule
$\tau$ & \normalfont\small R-REBEL $\ell_1$-std & \normalfont\small R-REBEL Huber-std & \normalfont\small GRPO + entropy & \normalfont\small GRPO, base-LM ref. \\
\midrule
$0$ & Awesome!!{\textquotedbl} gorgeous!!{\textquotedbl} appreciation \normalfont\scriptsize(100\%) & Awesome!!{\textquotedbl} thoughtful!!{\textquotedbl} loving \normalfont\scriptsize(76\%) & Review favourite Favorite favourite Favorite \normalfont\scriptsize(94\%) & We love it!!{\textquotedbl} beautiful \normalfont\scriptsize(59\%) \\[2pt]
$30$ & Awesome!!{\textquotedbl} gorgeous!!{\textquotedbl} appreciation \normalfont\scriptsize(56\%) & Awesome!!{\textquotedbl} gorgeous!!{\textquotedbl} adore \normalfont\scriptsize(83\%) & Review favourite Favorite favourite Favorite \normalfont\scriptsize(19\%) & We love it!!{\textquotedbl} beautiful \normalfont\scriptsize(49\%) \\[2pt]
$50$ & \textnormal{\scriptsize[U+304C]} marvelous Spanish wonderful Perfect \normalfont\scriptsize(47\%) & Excellent!!{\textquotedbl} delighted\allowbreak{}Interstitial Likes \normalfont\scriptsize(51\%) & Excellent favorite \textnormal{\scriptsize[U+203A]} Favorite \textnormal{\scriptsize[U+203A]} \normalfont\scriptsize(41\%) & {\textbackslash}n{\textbackslash}n{\textbackslash}n{\textbackslash}n{\textbackslash}n \normalfont\scriptsize(16\%) \\[2pt]
$70$ & \textnormal{\scriptsize[U+304C]} marvelous Portuguese Hindi wonderful \normalfont\scriptsize(44\%) & superb!!{\textquotedbl} impressed\allowbreak{}Interstitial Likes \normalfont\scriptsize(78\%) & Excellent Favorite \textnormal{\scriptsize[U+203A]} large\allowbreak{}Download Favorite \normalfont\scriptsize(24\%) & \textnormal{[RLO]}{\textbackslash}n{\textbackslash}n{\textbackslash}n \normalfont\scriptsize(53\%) \\[2pt]
$90$ & \textnormal{\scriptsize[U+00F0]} Catalan Catalan Catalan\textnormal{\scriptsize[U+30E5]} \normalfont\scriptsize(72\%) & Images\allowbreak{}Interstitial Likes $\times$ congratulations \normalfont\scriptsize(53\%) & ategories \textnormal{\scriptsize[U+203A]} large\allowbreak{}Download Favorite \textnormal{\scriptsize[U+203A]} \normalfont\scriptsize(73\%) & \textnormal{[RLO]}{\textbackslash}n{\textbackslash}n{\textbackslash}n \normalfont\scriptsize(91\%) \\[2pt]
\bottomrule
\end{tabular}

\caption{The most frequent prompt each Reward~II policy emits at five floors,
across the $\wNsent$ test sentences, with the share of sentences that receive
it (prompts are sampled). Decoded with the policy's tokenizer and shown
verbatim. \texttt{[RLO]} is an invisible right-to-left-override character
(U+202E), \texttt{\textbackslash n} a newline, and \texttt{[U+\dots]} other
non-ASCII characters by code point (e.g.\ U+304C is the Japanese kana
\emph{ga}).}
\label{tab:prompts}
\end{table}

\paragraph{A graded family of prompts.}
Reward~I's policies learned a few prompt modes (Section~\ref{sec:res-diagnosis}).
Under Reward~II, every content level has its own target, and the policies
learned families of related prompts whose expected content rises in steps.
Huber-std's family runs from effusive word lists at low floors
(\texttt{Awesome!!" gorgeous!!" adore}) through \texttt{Excellent!!"
delighted\dots} to \texttt{superb!!" impressed\dots}
(Table~\ref{tab:prompts}). Across all floors and test sentences it uses
$\wHuberPromptsTotal$ distinct prompts, against $\dgHuberPromptsTotal$ under
Reward~I. $\ell_1$-std found a different family, built around words for
languages (\texttt{marvelous Spanish}, \texttt{Portuguese}, \texttt{Catalan}).

\subsection{What the expectation floor costs: rewrites become bimodal}
\label{sec:res-cost}

\begin{table}[t]
\centering
\setlength{\tabcolsep}{3pt}
% generated by analysis/example_table.py -- do not edit by hand
\begin{tabular}{r>{\raggedright\arraybackslash\ttfamily\scriptsize}p{0.19\textwidth}cc>{\raggedright\arraybackslash\scriptsize}p{0.255\textwidth}>{\raggedright\arraybackslash\scriptsize}p{0.255\textwidth}}
\toprule
\multicolumn{6}{l}{\normalfont\small Input: \emph{do not go here if you are interested in eating good food.}} \\
\midrule
$\tau$ & \normalfont\small Prompt & $\mathbb{E}[c]$ / $\mathbb{E}[s]$ & Positive & A positive rewrite & A non-positive rewrite \\
\midrule
\multicolumn{6}{l}{\normalfont\small\emph{Reward~II (floor on expected content)}} \\[1pt]
$0$ & Awesome!!{\textquotedbl} thoughtful!!{\textquotedbl} loving & $27$ / $93$ & $94\%$ & i have been eating at this place for years and i am very happy with the food i have been getting. {\normalfont\tiny($c{=}26$, $s{=}100$)} & i just went there and the food is not what i expected and they did not have any chicken. {\normalfont\tiny($c{=}25$, $s{=}0$)} \\
$30$ & Awesome!!{\textquotedbl} gorgeous!!{\textquotedbl} adore & $26$ / $98$ & $98\%$ & very delicious! {\normalfont\tiny($c{=}27$, $s{=}100$)} & i like to go to the food hall for the lunch buffet. {\normalfont\tiny($c{=}32$, $s{=}1$)} \\
$50$ & Excellent!!{\textquotedbl} delighted\allowbreak{}Interstitial Likes & $60$ / $42$ & $42\%$ & good place {\normalfont\tiny($c{=}28$, $s{=}100$)} & do not go here if you are interested in eating good food. {\normalfont\tiny($c{=}100$, $s{=}0$)} \\
$70$ & superb!!{\textquotedbl} impressed\allowbreak{}Interstitial Likes & $77$ / $28$ & $28\%$ & very good! {\normalfont\tiny($c{=}23$, $s{=}100$)} & do not go here if you are interested in eating good food. {\normalfont\tiny($c{=}100$, $s{=}0$)} \\
$90$ & Images\allowbreak{}Interstitial Likes $\times$ congratulations & $89$ / $11$ & $12\%$ & congratulations {\normalfont\tiny($c{=}12$, $s{=}95$)} & do not go here if you are interested in eating good food. {\normalfont\tiny($c{=}100$, $s{=}0$)} \\
\midrule
\multicolumn{6}{l}{\normalfont\small\emph{Reward~I (floor on every sample)}} \\[1pt]
$0$ & Awesome!!{\textquotedbl} wonderful.{\textquotedbl} loving & $28$ / $95$ & $96\%$ & i've been to this place twice and i love the food. {\normalfont\tiny($c{=}28$, $s{=}100$)} & i am not a fan of the food, but it's not the best, i think. {\normalfont\tiny($c{=}30$, $s{=}1$)} \\
$30$ & Positive praise (\% contrasts ($-$ & $62$ / $77$ & $78\%$ & if you are interested in eating good food {\normalfont\tiny($c{=}65$, $s{=}100$)} & go here if you are not interested in eating good food. {\normalfont\tiny($c{=}72$, $s{=}5$)} \\
$50$ & Positive praise (\% contrasts ($-$ & $64$ / $76$ & $80\%$ & go here if you are looking for good food. {\normalfont\tiny($c{=}68$, $s{=}100$)} & do not go here if you are not interested in eating good food. {\normalfont\tiny($c{=}89$, $s{=}0$)} \\
$70$ & Sample conditional contrasts ($-$ contrasts & $71$ / $53$ & $52\%$ & if you are interested in eating good food, go here. {\normalfont\tiny($c{=}70$, $s{=}100$)} & you would go here if you were interested in eating good food. {\normalfont\tiny($c{=}78$, $s{=}44$)} \\
$90$ & Example $\times$ grep Example $\times$ & $90$ / $6$ & $6\%$ & do you want to go to {\normalfont\tiny($c{=}34$, $s{=}88$)} & do not go here if you are interested in eating good food. {\normalfont\tiny($c{=}100$, $s{=}0$)} \\
\bottomrule
\end{tabular}

\caption{\textbf{One input, five floors, both rewards} (\rrebel{} Huber-std,
step $12{,}000$, test sentence $\exSrcIndex$). For each floor: the policy's most
likely prompt, the expected content and sentiment of its rewrites ($50$
task-model samples), the share of rewrites that are positive (sentiment
$\ge 50$), and one rewrite from each group: the one with the group's median
content, not the best one. Rewrites are lower-cased by the pipeline. Under
Reward~II the middle floors are met by a mixture of exact copies and short
positive sentences. Under Reward~I the same floors produce rewrites that
change the sentiment and keep the content.}
\label{tab:example}
\end{table}

\begin{table}[t]
\centering
\small
% generated by analysis/example_table.py -- do not edit by hand
\begin{tabular}{lcccc}
\toprule
 & $\tau = 30$ & $\tau = 50$ & $\tau = 70$ & Success \\
 & \multicolumn{3}{c}{success / drift / copy (\% of rewrites)} & (mean) \\
\midrule
Reward~II, R-REBEL Huber-std & $28$ / $60$ / $1$ & $4$ / $53$ / $27$ & $2$ / $31$ / $54$ & $11$ \\
Reward~I, R-REBEL Huber-std & $35$ / $24$ / $5$ & $19$ / $25$ / $8$ & $9$ / $20$ / $10$ & $21$ \\
Reward~II, R-REBEL $\ell_1$-std & $28$ / $58$ / $2$ & $6$ / $54$ / $17$ & $2$ / $28$ / $43$ & $12$ \\
Reward~I, R-REBEL $\ell_1$-std & $34$ / $15$ / $6$ & $5$ / $3$ / $41$ & $3$ / $5$ / $41$ & $14$ \\
\bottomrule
\end{tabular}

\caption{\textbf{What individual rewrites look like at the binding floors}
($\exAuditN$ test sentences, every fifth one, $50$ rewrites each; prompts
decoded as in each campaign's evaluation). \emph{Success}: positive (sentiment
$\ge 50$) \emph{and} at or above the floor. \emph{Drift}: positive but below the
floor. \emph{Copy}: not positive, with content $\ge 95$, i.e.\ the input
reproduced. The last column averages success over the three floors.}
\label{tab:outcomes}
\end{table}

Reward~II asks for the right \emph{expected} content and as much
\emph{expected} sentiment as possible. A rewrite that is both positive and
faithful to the source is one way to deliver that. Another way, which the
expected values cannot tell apart from the first, is to sometimes copy the
source verbatim, which buys content, and sometimes write a short positive
sentence, which buys sentiment. Proposition~\ref{prop:eps} already showed that
the expectation constraint is met by mixing. The trained policies do the mixing
not by randomizing between prompts but inside each prompt, through the task
model. Table~\ref{tab:example} shows one input. Under Reward~II, at $\tau = 50$
and above, the non-positive rewrites are the input copied verbatim and the
positive ones are generic (\emph{good place}, \emph{very good!}). Under
Reward~I, the same policy architecture produces rewrites that keep the content
and flip the sentiment.

Table~\ref{tab:outcomes} measures this over $\exAuditN$ test sentences. For
Huber-std at $\tau = 50$, $\exFiveHuberSuccessAtFifty\%$ of Reward~II rewrites are
positive and meet the floor, against $\exFourHuberSuccessAtFifty\%$ under
Reward~I. $\exFiveHuberCopyAtFifty\%$ are copies and $\exFiveHuberDriftAtFifty\%$
positive rewrites that fall below the floor. At $\tau = 70$, copies are
$\exFiveHuberCopyAtSeventy\%$ of Reward~II rewrites. Averaged over the three floors,
Huber-std's success rate halves, from $\exFourHuberSuccess\%$ to
$\exFiveHuberSuccess\%$. For $\ell_1$-std the rates are similar under the two
rewards ($\exFourLoneSuccess\%$ and $\exFiveLoneSuccess\%$), because under
Reward~I its copy prompt had already taken over the middle floors
(Section~\ref{sec:res-diagnosis}). The per-sample reward column of
Table~\ref{tab:full}, computed on all $\wNsent$ sentences, shows the same
direction: $\wHuberRew$ for Huber-std under Reward~II against $\vHuberRew$
under Reward~I.

The conclusion is not that Reward~II is mis-specified. It optimizes exactly the
curve of Eq.~\eqref{eq:curve}, and it does so better than Reward~I does. The
conclusion is that \emph{an expected-value tradeoff curve is the wrong target
when every individual rewrite should be a faithful transfer}. For that goal,
the constraint has to act on individual rewrites, as Reward~I's does.
Reward~I's own defects then have to be fixed separately: the right curve plots
the floor against sentiment, and the policy has to stop clustering. We return to
this in Section~\ref{sec:limitations}.

\subsection{Learning dynamics on the test split}
\label{sec:res-dynamics}

\begin{figure}[t]
\centering
\IfFileExists{figures/test_learning_v5.pdf}{\includegraphics[width=0.9\textwidth]{test_learning_v5.pdf}}{\fbox{\parbox{0.8\textwidth}{\centering test-split learning curves pending}}}
\caption{\textbf{Test-split learning curves} (Reward~II). Each arm's
checkpoints at \wLcNsteps{} matched steps, evaluated with the full protocol (one
evaluation seed; at step $12{,}000$ Huber-std's hypervolume varies across five
seeds with standard deviation \wHuberHVseedSD). \textbf{Left:} hypervolume of
the tradeoff curve. \textbf{Right:} mean $|\mathbb{E}[c]-\tau|$ over the binding
floors.}
\label{fig:dynamics}
\end{figure}

Figure~\ref{fig:dynamics} evaluates every arm's checkpoints at \wLcNsteps{}
matched steps (\wLcFirst--\wLcLast). Huber-std's hypervolume exceeds
GRPO + entropy's at every checkpoint from step \wLcHVLeadFrom, by
\wLcHVLeadMin--\wLcHVLeadMax{} points. Both \rrebel{} arms improve throughout
training: Huber-std from \wHuberHVfirst{} to \wHuberHV, and $\ell_1$-std from
\wLoneHVfirst{} to \wLoneHV. GRPO + entropy stays within
\wGentHVLcMin--\wGentHVLcMax{}, and the base-LM-reference arm within
\wGbaseHVLcMin--\wGbaseHVLcMax. Huber-std tracks the floor closely from the first checkpoint on (mean
deviation \wHuberCalibFirst{} at step \wLcFirst, between \wHuberCalibLcMin{}
and \wHuberCalibLcMax{} thereafter). The GRPO arms start further from it
(\wGentCalibFirst{} and \wGbaseCalibFirst) and then fluctuate between
\wGentCalibLcMin--\wGentCalibLcMax{} and
\wGbaseCalibLcMin--\wGbaseCalibLcMax{} respectively. Under Reward~I,
$\ell_1$-std regressed late in training. Under Reward~II it does not, and the
mechanism of that regression (Section~\ref{sec:res-diagnosis}) explains why.

\subsection{Why the per-sample floor clusters (Reward~I)}
\label{sec:res-diagnosis}

\begin{figure}[!tp]
\centering
\includegraphics[width=0.86\textwidth]{frontier_panels.pdf}
\caption{\textbf{The same policies trained with Reward~I} (v4, step $12{,}000$),
labelled by floor. Many consecutive floors land on one point and share a range
label (e.g.\ $\tau = 50$--$95$ for $\ell_1$-std). The points also sit far from
$\mathbb{E}[c] = \tau$: every floor from $\tau = 50$ up gets $\ell_1$-std's copy
prompt, with content $\approx\dgLoneCopyC$. In v4 the evaluation used the legacy decoder, which samples from the
top three tokens and is near-greedy for these policies.}
\label{fig:frontier-v4}
\end{figure}

\begin{figure}[t]
\centering
\includegraphics[width=\textwidth]{menu_mechanism.pdf}
\caption{\textbf{Reward~I picks between a few prompts, and its switch points are
not at the floors} (v4, step $12{,}000$). \textbf{Top:} the Reward~I value of
each prompt the policy uses, against the floor, from per-sample re-evaluation on
all $\vNsent$ test sentences. Prompts within $3$ points of each other in
(content, sentiment) share a colour; dotted black: the best prompt. \textbf{Bottom:}
the policy's share of each prompt across the test sentences. Dotted vertical
lines mark where the reward-optimal prompt changes.}
\label{fig:menu}
\end{figure}

Figure~\ref{fig:frontier-v4} shows the Reward~I curves of v4, labelled by floor.
Two things are wrong with them. Consecutive floors coincide, and points sit far
from $\mathbb{E}[c] = \tau$. To find out why, we re-scored the
$\dgMenuSize$ prompts that account for at least $98\%$ of the four policies'
test-split choices on all $\vNsent$ test sentences, keeping every task-model
sample. We also mapped each policy's most likely prompt on a $200$-point grid
of $\lmb$.

\paragraph{Each policy is a lookup table of a few prompts.}
$\dgLoneCoverK$ prompts cover $99\%$ of $\ell_1$-std's $10{,}000$
(sentence, floor) choices, $\dgHuberCoverK$ cover Huber-std's and
$\dgGentCoverK$ GRPO + entropy's. Along the fine $\lmb$ grid the most likely
prompt changes only $\dgLoneSwitches$, $\dgHuberSwitches$ and $\dgGentSwitches$
times respectively. The policies are also nearly input-independent: at a given
floor, almost every sentence receives the same prompt. Lemma~\ref{lem:cluster}
then bounds the curve to a handful of distinct points, which is the clustering
of Figure~\ref{fig:frontier-v4}. Choosing prompts per sentence would avoid it:
re-solving Reward~I per sentence over the same $\dgMenuSize$ prompts gives
$\dgDesignSampleX$ distinct points, against $\dgDesignSampleOne$ for the best
single prompt per floor. Of the constraint forms we compared, only the floor on
the \emph{policy's} expected content removes the clustering without
per-sentence choice, giving $\dgDesignMixOne$ distinct points, because its
optimum mixes prompts with weights that vary continuously with $\tau$
(Proposition~\ref{prop:eps}).

\paragraph{The switch points are set by feasibility, not by the floor.}
The policies switch prompts close to where Reward~I says they should: against
the best of its own prompts at each floor, Huber-std and GRPO + entropy give up
at most $\dgRegretGoodMax$ reward points (Figure~\ref{fig:menu}). The problem
is where Reward~I puts the switches. The reward is sentiment times the
probability of meeting the floor, and a prompt's per-sample content is widely
spread: its standard deviation within one sentence is $\dgWithinSDmin$--$\dgWithinSDmax$
points, and copy prompts reproduce the source exactly in only
$\dgCopyExactMin$--$\dgCopyExactMax\%$ of samples. So the reward keeps the most
positive prompt until only $\dgFirstOutLo$--$\dgFirstOutHi\%$ of its samples
meet the floor, then jumps to a prompt that meets it
$\dgFirstInLo$--$\dgFirstInHi\%$ of the time. Just before a switch, $\mathbb{E}[c]$
lies below $\tau$. Just after, it can lie up to $\vGapAllMax$ points above. Within each
plateau $\mathbb{E}[c]$ is constant while $\tau$ rises, which produces the
sawtooth of Figure~\ref{fig:compare} (left).

\paragraph{The late regression of $\ell_1$-std under Reward~I.}
In v4, $\ell_1$-std's test reward peaked at step \vLonePeakStep{}
(\vLonePeakRew) and fell by \vLoneDrop{} points by step $12{,}000$. The
mechanism is visible in the prompt assignment. The lowest floor at which its copy
prompt (content $\approx\dgLoneCopyC$, sentiment $\approx\dgLoneCopyS$) serves most sentences fell
from $\tau = \dgLoneTakeoverPeak$ at the peak to $\tau = \dgLoneTakeoverEnd$ at
the end, so the reward at $\tau = 50$ fell from $\dgLoneRewFiftyPeak$ to
$\dgLoneRewFiftyEnd$. The floors $\tau = 40$--$70$ account for
$\dgLoneDropMid$ of the $\dgLoneDropAll$-point drop. At those floors its own mid
prompt (content $\approx\dgLoneMidC$, sentiment $\approx\dgLoneMidS$) would have earned up to
$\dgLoneRegretMax$ more reward points. Under Reward~I, a floor the copy prompt
over-satisfies costs nothing in the constraint term, so nothing
holds the copy prompt back once it has spread. Under Reward~II the over-satisfied
floors pay for the lost sentiment, and the regression does not occur
(Section~\ref{sec:res-dynamics}). Development-set early stopping also recovers
it under Reward~I: it selects step \vLoneDevStep, with test reward
\vLoneDevRew.

\paragraph{The evaluation decoder.}
The code's ``greedy'' decoder, used in every evaluation before v5, in fact
samples each token from the top three. For the three policies with peaked
distributions this returns the most likely prompt with probability
$\dgHuberPArgmax$ (median over sentences and floors), and re-scoring with the true
most likely prompt changes their Reward~I reward by at most
$\dgPeakedArgmaxMaxAbs$ points. For the base-LM-reference arm, which is less
peaked, the change is $\dgGbaseArgmaxRewDiff$. None of the v4 comparisons
depend on this. Because the sampler is seeded, prompts were identical across
evaluation seeds on the same GPU model. Across GPU models the random stream
differs, and the fraction of prompts that changed matches what two independent
draws would give: $\dgHuberFlipExpected\%$ expected against
$\vHuberFlipMax\%$ observed for Huber-std, and $\dgGbaseFlipExpected\%$ against
$\vGbaseFlipMax\%$ for the base-LM-reference arm.

\subsection{How much the reward scale varies across contexts}
\label{sec:res-scale}

\begin{figure}[t]
\centering
\includegraphics[width=\textwidth]{reward_scale.pdf}
\caption{\textbf{Reward scale varies by orders of magnitude across contexts}
(Reward~I; \rrebel{} $\ell_1$-std at step \sProbeFinal; $40$ test sources
$\times$ $10$ floors, $G = 16$ prompts per group drawn by the training sampler).
\textbf{Left:} in-group reward standard deviation $\hat\sigma(x,\lmb)$ by floor
(boxes: interquartile range; whiskers: 5th--95th percentile).
\textbf{Right:} the one-window KL radius that Proposition~\ref{prop:trust}
assigns to each group, computed with the group's own empirical distribution as
the reference (so it is capped at $\log G = \sLogG$ nats, a point mass). Under
the plain target, one global $\beta$ leaves some contexts unmoved and
collapses others; under the standardized target, the radii concentrate.}
\label{fig:scale}
\end{figure}

Section~\ref{sec:th-std} argued that standardization matters when the
reward's scale varies across contexts. To measure that variation directly, we
drew groups exactly as a training step does ($G = 16$ prompts per
$(x,\lmb)$, top-$k=200$ sampling, $N=50$ task-model samples per prompt), from
the v4 \rrebel{} $\ell_1$-std checkpoints at steps \sProbeSteps, for $40$ test
sources and $10$ floors each. The measurement uses Reward~I. Reward~II's
in-group spread varies across contexts for an analogous reason: its content
weight $\mu_g$ is $0$ where a group meets its floor and grows with the
shortfall where it does not. Figure~\ref{fig:scale} shows the
result.

\paragraph{The scale spans two to three orders of magnitude.}
Across groups, the in-group standard deviation $\hat\sigma$ varies by a factor
of \sRangeQ{} between its 5th and 95th percentiles. The variation is
systematic in $\lmb$. At the final probed checkpoint the median $\hat\sigma$ is
\sSigZeroFin{} at $\tau = 0$, peaks at \sSigMaxFin{} for low-to-middle floors,
and falls to \sSigNinetyFin{} at $\tau = 90$: a \sTauRatioFin-fold spread in
median alone. The mechanism is the two-branch reward. Where most samples meet
the floor, rewards are sentiment scores spread over tens of points. Where few
do, they are penalties of size $0.01\times$ the shortfall. The contexts with the
tightest floors, where few prompts are feasible, are exactly the ones with the
smallest reward differences.

\paragraph{One global temperature cannot serve all contexts.}
Applying Proposition~\ref{prop:trust} group by group quantifies the
consequence. With the plain target $\polref e^{r/\beta}$ and our
$\beta = 0.5$, \sPlainStuck{} of groups would move less than $0.01$ nats in a
window --- no learning --- while \sPlainCollapse{} would jump to within
$0.27$ nats of a point mass. Raising $\beta$ to rescue the second set freezes more
of the first, and lowering it does the reverse. With the standardized target,
\sStdStuck{} of groups fall below $0.01$ nats and \sStdCollapse{} come near the
point-mass cap; the median radius is \sStdMedKL{} nats. This is the practical
content of Theorem~\ref{thm:consistency}(ii) and Proposition~\ref{prop:trust}:
in a $\lmb$-conditioned problem, dividing by the in-group standard deviation is
what lets a single $\beta$ mean the same thing at every point on the curve.
Every \rrebel{} arm we train uses standardization, so this is a measurement of
the problem, not an ablation of the method.

\subsection{Training-seed replication (v3, three seeds, development set)}
\label{sec:seed-replication}

The v4 and v5 campaigns have one training run per arm, so they cannot measure how much a
result depends on the training seed. The earlier v3 campaign can, with a
caveat: it used the same code, hyperparameters and budget, but its only
surviving evaluation is the in-training score on a $10$-sentence development
set, over the $10$-point grid $\{0,\dots,0.9\}$ (Section~\ref{sec:provenance}).
Table~\ref{tab:main} reports it at step $11{,}500$, the largest step at which
all fifteen runs were evaluated.

\begin{table}[t]
\centering
\begin{tabular}{lrrrc}
\toprule
Algorithm & Reward & Content & Sentiment & Seeds \\
\midrule
R-REBEL ($\ell_1$-std) & 39.34 & 56.25 & 53.70 & 3/3 \\
R-REBEL (Huber-std) & 39.27 & 58.62 & 49.55 & 3/3 \\
R-REBEL ($\ell_1$-ent) & 37.49 & 58.17 & 48.15 & 3/3 \\
GRPO + entropy & 33.36 & 53.57 & 50.57 & 3/3 \\
GRPO & 30.19 & 55.92 & 43.03 & 3/3 \\
\bottomrule
\end{tabular}

\caption{\textbf{v3 campaign} (three training seeds per arm): mean reward,
content and sentiment at matched step $11{,}500$, on the \emph{$10$-sentence
development set} and the $10$-point $\lmb$ grid. Not comparable in absolute
terms to the v4 or v5 test-split numbers.}
\label{tab:main}
\end{table}

\paragraph{The family ordering holds across seeds.}
The best \rrebel{} variants score $39.34$ ($\ell_1$-std) and $39.27$
(Huber-std), tied within noise, against $33.36$ for GRPO + entropy, a gap of
$5.98$ points. There is no overlap at the level of individual runs: the worst
\rrebel{} seed ($36.82$) beats the best GRPO seed ($34.65$). GRPO's seeds are
also far more dispersed, spanning $26.3$--$34.7$ against $36.8$--$40.0$ for the
nine \rrebel{} seeds. A single GRPO run is therefore a less reliable estimate of
its arm than a single \rrebel{} run. $\ell_1$-ent ($37.49$) trails $\ell_1$-std
by $1.85$ points: adding an entropy bonus to a -std variant hurt, consistent
with Proposition~\ref{prop:mirror}, under which each window already moves by a
bounded KL. None of the three v3 $\ell_1$-std seeds shows the late decline of the
v4 $\ell_1$-std run (Section~\ref{sec:res-diagnosis}) on this development set.

\paragraph{The baseline was debugged, and plain GRPO is a collapse diagnostic.}
\label{sec:baseline-defense}
Three GRPO configurations compared at matched step $4500$ show the baseline
being repaired rather than tuned against:

\begin{center}
\begin{tabular}{lrl}
\toprule
GRPO configuration (v3, dev set) & Mean reward & Change \\
\midrule
Rolling reference anchor (defective) & $27.16$ & --- \\
Pinned reference anchor (Section~\ref{sec:grpo}) & $31.30$ & $+4.14$ \\
\quad + entropy bonus $\alpha_H = 0.01$ & $33.10$ & $+1.80$ \\
\midrule
\rrebel{} (all variants, same step) & $39.26$ & $+6.16$ remaining \\
\bottomrule
\end{tabular}
\end{center}

The rolling anchor re-copied the KL reference from the live policy every $150$
steps, which zeroed the penalty and its gradient at every refresh. Pinning it
was worth $+4.14$ points, and the entropy bonus a further $+1.80$. The
plain-GRPO arm, however, never learns. Its score moves from $30.38$ at step
$500$ to $30.19$ at step $11{,}500$. Its logged KL sits at $9.824$, against the
$k_3$ estimator's ceiling of $9.8249$, and it explores a median of $5$ distinct
prompt tokens per cycle: one frozen prompt. Because its pinned anchor is the
uniform distribution (Section~\ref{sec:algos-caveat}), plain GRPO and
GRPO + entropy are one configuration at two strengths of the same entropy
lever, not two independent baselines. This is why v4 and v5 replace plain GRPO with
the base-LM-reference arm.

\begin{figure}[t]
\centering
\includegraphics[width=0.78\textwidth]{learning_curves.pdf}
\caption{\textbf{v3 campaign}: mean reward on the $10$-sentence development set
against training step, averaged over three seeds, with the band spanning the
seed minimum to maximum. Blue: \rrebel{} variants; orange: GRPO arms.}
\label{fig:learning}
\end{figure}

\subsection{Data provenance}
\label{sec:provenance}

The v3 campaign ran to completion, but a purge of the scratch filesystem then
destroyed all of its checkpoints and per-$\lmb$ evaluation files. Its numbers
were recovered from the experiment tracker's scalar history, by a
reconstruction that is validated three ways per run and refuses to emit data if
any check fails. All fifteen runs pass. The tracker logged only
$\lmb$-averaged scores, so no v3 tradeoff curve can be drawn, and three runs lack
a step-$12{,}000$ evaluation (hence step $11{,}500$ above). The v4 campaign was
run to regenerate what was lost, and v5 with the same safeguards. Their
checkpoints are mirrored off the purgeable filesystem on every save, and every
evaluation file is archived. Every v4 and v5 number in this paper is regenerated
from those files by \texttt{analysis/full\_eval.py}.

